\documentclass[11pt]{article}
\usepackage[margin=0.85in]{geometry}
\usepackage{amsmath,amssymb,amsthm,booktabs,hyperref,microtype}
\setlength{\emergencystretch}{2em}
\newtheorem{theorem}{Theorem}
\newtheorem{lemma}{Lemma}
\newtheorem{proposition}{Proposition}
\newcommand{\E}{\mathbb E}
\newcommand{\Prb}{\mathbb P}
\newcommand{\Poi}{\operatorname{Poi}}
\newcommand{\eff}{\operatorname{eff}}
\newcommand{\SmallUpperDisplay}{0.331}
\newcommand{\BridgeUpperDisplay}{0.301518}

\title{A Sufficient-Constant Proof for the Unchanged Support Estimator\\$C_0=70$, $L_0=1$}
\author{Independent analytic derivation and reproducible interval certificates}
\date{November 4, 2026}
\begin{document}
\maketitle
\begin{abstract}
We prove the original single-run, multiscale support guarantee with real
Poisson mean $70n/[\epsilon(1+\log n)]$ for every positive integer $n$.
The new part is uniform control of the upper tail below the previously
certified threshold $1+\log n=14$. It combines a Poisson occupancy
homogenization lemma, a large-mean anchor category, and exact moment
generating function certificates. All degree ceilings and scale transitions
are covered. Simulations are not premises of this proof. This is a
sufficiency result for 70, not a proof of the optimal real constant.
\end{abstract}

\section{Exact algorithm and theorem}
Let $p$ be a probability distribution on a finite or countable alphabet.
For $n\ge1$ and $0<\epsilon<1$, put
\[
 L=1+\log n,\qquad m=\frac{70n}{\epsilon L},\qquad
 N_i\ \hbox{independent},\quad N_i\sim\Poi(mp_i),\qquad
 D=\sum_i1_{N_i>0}.
\]
The mean $m$ is a real number, not an integer ceiling. Write
$S=|\operatorname{supp}p|$,
$\Delta_r=\sum_{i>r}p_{(i)}$, and
$e_\epsilon=\min\{r\ge1:\Delta_r\le\epsilon\}$.
The latter is finite even if $S=\infty$.
Set
\[
 J=\left\lfloor\log_2\frac{256L^2}{70}\right\rfloor,\quad
 A_j=\frac{70\,2^j}{8L},\quad B=64L,\quad
 \delta_j=2^{-j-2},\quad
 k_j=\left\lceil\sqrt{B/A_j}(j+3)\log2\right\rceil.
\]
Let $a_j=A_j/m$, $b=B/m$, and
\[
 P_j(x)=-\frac{T_{k_j}((b+a_j-2x)/(b-a_j))}
                  {T_{k_j}((b+a_j)/(b-a_j))}
       =\sum_{\ell=0}^{k_j}c_{j,\ell}x^\ell.
\]
Use $f_j(0)=0$, $f_j(\ell)=1+c_{j,\ell}\ell!/m^\ell$ for
$1\le\ell\le k_j$, and $f_j(\ell)=1$ for $\ell>k_j$. Define
\[
 Z_j=\sum_i f_j(N_i),\qquad
 \widetilde S=\max\left\{D,\max_{0\le j\le J}
                   \frac{Z_j-\delta_j n}{1+\delta_j}\right\},\quad
 \widehat S=\min\{n,\max\{1,\lceil\widetilde S\rceil\}\}.
\]
The binary output is Yes if and only if $\widetilde S\le n$.
These are the original scales, degree ceilings, correction, and output;
there is no degree cap or modification of the statistic in the theorem.
For countable support, $\sum_iN_i\sim\Poi(m)<\infty$ almost surely.
Since $f_j(0)=0$, only finitely many terms in each $Z_j$ are nonzero
almost surely; the statistic is therefore well-defined.

\begin{theorem}\label{thm:main}
For every such $n,\epsilon,p$, the failure probability is less than
$0.331$, and in particular
\[
 \Prb\{\min(e_\epsilon,n)\le\widehat S\le\min(S,n)\}>\frac23.
\]
On the promise $S\le n$, the binary output is Yes; on the promise
$\Delta_n\ge\epsilon$, it is No, with probability greater than $2/3$.
On either promise, the binary output and the displayed interval are
simultaneously correct with probability greater than $2/3$.
There is no binary correctness claim in the intervening gap.
\end{theorem}

The upper and lower deviations are different events. Establishing a bound
for each at $1/3$ would not prove the simultaneous statement. Below we
explicitly add the bounds where necessary.

\section{Well-defined parameters and exact coefficient formula}
Since $256/70>1$, $J\ge0$ for all $L\ge1$, and its definition gives
$0<A_j\le B/2$. With
$\alpha_0=3\log2\sqrt{512/70}$, the sequence
$(j+3)2^{-j/2}$ is decreasing, and hence
\[
 k_j\le\alpha_0L+1<6.7L<64L/\sqrt2
 \le\sqrt{B(B-A_j)}.
\]
The usual Chebyshev normalization bound gives
$\rho_j=T_{k_j}((B+A_j)/(B-A_j))^{-1}\le\delta_j^2/2$.
The constant inequalities in this paragraph are checked by outward
intervals as well as the finite implementation audit.

The original expectation lemma was stated with $b\le1$. Its proof is
valid without that restriction: the intervals $[0,a]$, $[a,b]$, and
$[b,\infty)$ use respectively concavity, $|T_k|\le1$, and the inequality
$m\sqrt{b(b-a)}\ge k$. None uses an absolute upper bound on $b$.
Thus removing the implementation's old $b\le1$ guard does not change
the estimator and does not invalidate that lemma. The new finite-range
upper proof below in fact uses the exact coefficients directly.

For $k\ge1$,
\[
 T_k(1+z)=\sum_{r=0}^k d_{k,r}z^r,\qquad
 d_{k,r}=\frac{k\,2^r(k+r-1)!}{(k-r)!(2r)!}>0.
\]
This identity follows either by differentiating the Chebyshev differential
equation at 1 or by its polynomial recurrence. Put
$z=2A/(B-A)$ and $\beta=2/(B-A)$. Then
\begin{equation}\label{eq:coeff}
 u_\ell=
 \frac{\beta^\ell\ell!\sum_{r=\ell}^k d_{k,r}\binom r\ell z^{r-\ell}}
      {\sum_{r=0}^k d_{k,r}z^r}>0,\qquad
 f(\ell)-1=(-1)^{\ell+1}u_\ell.
\end{equation}
All terms used to enclose $u_\ell$ are positive. The independent checker
also uses
$d_{k,r+1}/d_{k,r}=(k^2-r^2)/[(r+1)(2r+1)]$ and a different polynomial
evaluation to avoid relying on just one coefficient implementation.

\section{Lower deviation on the new range}
\begin{lemma}[Discovery coupling]\label{lem:discovery}
If $\Delta_r\ge\epsilon$ and $r\le n$, then
\[
 \Prb(D\le r)\le\Prb\{\Poi(m\epsilon)\le r\}
                 \le\Prb\{\Poi(70n/L)\le n\}.
\]
\end{lemma}
\begin{proof}
Until $r+1$ distinct labels have been seen, the probability that the next
observation is new is at least $\epsilon$, since the total mass of any
set of at most $r$ labels is at most $1-\Delta_r$.
Thin the new-label event, using an extra independent uniform variable,
to have conditional probability exactly $\epsilon$. After $r+1$
discoveries, continue with independent Bernoulli($\epsilon$) marks.
These marks are iid and are a thinning of the Poisson arrival process;
before the target is reached each mark forces a discovery.
Thus $D\le r$ implies that at most $r$ marks arrived. Poisson thinning
and monotonicity in the threshold give the assertion.
\end{proof}
For $C\ge69$ and $1\le L\le32$, Chernoff gives
\[
 \Prb\{\Poi(Cn/L)\le n\}
 \le\exp\{-e^{L-1}[C/L-1-\log(C/L)]\}.
\]
On each of the 992 closed intervals $[a,b]$ of width $1/32$ covering
$[1,32]$, replace $e^{L-1}$ by $e^{a-1}$ and $C/L$ by $69/b$.
The function $x-1-\log x$ is increasing for $x>1$.
Outward interval evaluation proves the uniform bound
\begin{equation}\label{eq:lower}
 R_{\rm small}<1.657\cdot10^{-27}.
\end{equation}
This is much smaller than the upper-tail margins used below. For the
interval output take $r=\min(e_\epsilon,n)-1$ when it is positive;
minimality of $e_\epsilon$ gives $\Delta_r>\epsilon$.
For the binary alternative take $r=n$.

\section{Exact MGF and its finite certificates}
For a fixed scale write $Y=f(\Poi(\lambda))-1$. Set
$y_0=-1$, $y_\ell=(-1)^{\ell+1}u_\ell$ for $1\le\ell\le k$.
The other values of $Y$ are zero. Therefore, exactly,
\begin{equation}\label{eq:mgf}
 G(t,\lambda)=\E e^{tY}
 =1+e^{-\lambda}\sum_{\ell=0}^k(e^{ty_\ell}-1)
                                  \frac{\lambda^\ell}{\ell!}.
\end{equation}
In particular the zero-count contribution is negative and is retained.
If $G(t,\lambda)\le1+\eta$ uniformly for $\lambda\ge0$, exponential
Markov and independence imply
\begin{equation}\label{eq:chernoff}
 \Prb\{Z_j>(1+\delta_j)S+\delta_jn\}
 \le\exp\{S\log(1+\eta)-t\delta_j(n+S)\}.
\end{equation}
For fixed $n$, its logarithm is affine in $S$; the maximum over
$1\le S\le n$ is attained at $S=1$ or $S=n$.

\begin{lemma}[Polynomial certificate]\label{lem:polynomial}
Let $d\ge k$, $t,\eta>0$. If
\[
 Q(\lambda)=\eta\sum_{\ell=0}^d\frac{\lambda^\ell}{\ell!}
       -\sum_{\ell=0}^k(e^{ty_\ell}-1)\frac{\lambda^\ell}{\ell!}
 \ge0\qquad(\lambda\ge0),
\]
then $G(t,\lambda)\le1+\eta$ for every $\lambda\ge0$.
\end{lemma}
\begin{proof}
The Taylor polynomial is a lower bound for $e^\lambda$ on this domain.
Consequently $\eta e^\lambda$ is no smaller than the second sum;
multiply by $e^{-\lambda}$ and use \eqref{eq:mgf}.
\end{proof}
There are two separate implementations of the global nonnegativity test.
First, map a finite interval $[a,b]$ to $[0,1]$ and express $Q$ in the
Bernstein basis $\binom di x^i(1-x)^{d-i}$. Nonnegative coefficients
certify nonnegativity on the entire interval. Subdivide intervals if
needed. For $\lambda\ge R$, nonnegative coefficients of $Q(R+x)$
suffice. Second, for $d=k$ map the whole half-line by
$\lambda=x/(1-x)$; the Bernstein coefficients of
$(1-x)^k Q(x/(1-x))$ are explicitly
\[
 (\eta-(e^{ty_\ell}-1))\frac{(k-\ell)!}{k!}.
\]
de Casteljau subdivision proves nonnegativity on $[0,1]$ and hence on
the complete original half-line. The endpoint $x=1$ causes no gap:
all finite $\lambda$ correspond to $x<1$.
All these operations use outward intervals. A failed or ambiguous
positivity test is rejected; it is never replaced by a grid estimate.

\subsection{Integers $74\le n\le403$}
Every scale of these integers is included. The 329 integers from 75 to
403 give 2570 scale obligations; a separate bridge certificate covers
all seven scales at 74. A nonrigorous numerical search proposes rational
$t,\eta$. Lemma~\ref{lem:polynomial} then verifies every proposal over
the whole half-line. The final proof uses only the rational proposals
and the successful certificates, not correctness of the optimizer.
Summing \eqref{eq:chernoff}, with both $S$ endpoints checked, gives
\begin{equation}\label{eq:finite}
 U_{75:403}<0.291304,\qquad U_{74}<\BridgeUpperDisplay.
\end{equation}
The exact certified maximum is recorded in the accompanying JSON.
The independent replay rejects missing integers, missing scales,
incorrect degrees or corrections, and nonpositive rational parameters.
Deleting an integer or a scale is included as a negative control.

\subsection{Continuous range $7\le L\le14$}
Partition this interval into 700 closed rational slabs of width $1/100$.
For a slab $[a,b]$ put $N=\lfloor e^{a-1}\rfloor\le n$.
At scale $j$ choose
\[
 t=\frac4{\delta_jN},\qquad\eta=\frac4N=t\delta_j.
\]
For each possibly active scale include every integer between
$\lceil\alpha_ja\rceil$ and $\lceil\alpha_jb\rceil$, where
$\alpha_j=\sqrt{512/(70\,2^j)}(j+3)\log2$.
Equation~\eqref{eq:coeff} encloses the coefficients over the whole slab,
even when the ceiling or $J$ changes inside it.
Both nonnegativity implementations certify all 6597 resulting $(j,k)$
obligations. Thus $G\le1+\eta$. Using $\log(1+\eta)\le\eta$ in
\eqref{eq:chernoff} gives
\[
 S\log(1+\eta)-t\delta_j(n+S)\le-4n/N\le-4.
\]
There are at most 10 active scales for $L\le14$, so
\begin{equation}\label{eq:mid}
 U_{7:14}\le10e^{-4}<0.183157.
\end{equation}
This is continuous coverage, not a finite sample of $L$ values.

\section{A homogenization theorem for zero and singleton counts}
\begin{lemma}\label{lem:homog}
Let $X_i$ be independent $\Poi(\lambda_i)$, $i=1,\ldots,r$,
$F_0=\#\{i:X_i=0\}$, $F_1=\#\{i:X_i=1\}$.
For $0\le y\le1$ and nonnegative nondecreasing
$\phi:\{0,\ldots,r\}\to\mathbb R$,
\[
 \sup_{\lambda_1,\ldots,\lambda_r\ge0}\E[y^{F_0}\phi(F_1)]
 =\sup_{\lambda\ge0}\E_{\lambda_i=\lambda}[y^{F_0}\phi(F_1)].
\]
For computing the supremum one may restrict to $\lambda\ge1$.
\end{lemma}
\begin{proof}
First let $0<d=1-y\le1$ and all parameters be finite and positive.
For a single coordinate put
\[
 b=\lambda e^{-\lambda},\quad t=1-de^{-\lambda},\quad
 h=t-b=1-(d+\lambda)e^{-\lambda},\quad x=\log t.
\]
The unnormalized singleton generating polynomial after weighting zeros
by $y$ is $h+bz$. We show that $\log b$ is concave and $\log h$ convex
as functions of $x$. Since $dx/d\lambda>0$, for the first claim it
suffices that
\[
 \frac{d}{d\lambda}\frac{d\log b}{dx}
 =\frac{d-e^\lambda(\lambda^2-\lambda+1)}{d\lambda^2}\le0.
\]
Indeed $e^\lambda(\lambda^2-\lambda+1)\ge1$, as its derivative is
$e^\lambda\lambda(\lambda+1)\ge0$ and its value at zero is 1.
For the second claim,
\[
 \frac{d\log h}{dx}=\frac{t(\lambda-y)}{dh},\qquad
 \frac{d}{d\lambda}\frac{d\log h}{dx}
 =\frac{W(\lambda,d)}{de^{2\lambda}h^2},
\]
where
\[
 W=e^{2\lambda}-e^\lambda[\lambda^2-\lambda+1+d(\lambda+1)]+d.
\]
This decreases with $d$, and at $d=1$ equals
$e^\lambda(2\cosh\lambda-\lambda^2-2)\ge0$, by the nonnegative
remaining terms in the cosh power series.

Replace two $x$ coordinates by their arithmetic mean, with resulting
common parameter denoted by a star. Then
\[
 t_*^2=t_1t_2,\qquad b_*^2\ge b_1b_2,\qquad h_*^2\le h_1h_2.
\]
For two singleton indicators the total unnormalized mass, tail mass at
least 1, and tail mass at least 2 are respectively
$t_1t_2$, $t_1t_2-h_1h_2$, $b_1b_2$.
The replacement preserves the first and increases the latter two.
Every nondecreasing $\phi$ is a nonnegative linear combination of a
constant and upper-tail indicators. Conditional on the other coordinates,
the weighted expectation therefore does not decrease.

Repeatedly average the largest and smallest $x$ coordinates. Their mean
is unchanged; the sum of squared deviations decreases by half their
squared difference and consequently converges to zero. Continuity shows
that the fully equal vector has at least the original expectation.
The reverse supremum inequality is immediate.
Boundary points follow by continuity. If $y=0$ and a mean is zero,
the weighted expectation is zero. Infinite means are limits.
If $y=1$, the $F_1$ summands are independent Bernoulli variables with
probabilities $\lambda_i e^{-\lambda_i}$, each maximized at 1;
monotonicity of $\phi$ proves this remaining case.

Finally, conditional on the other coordinates the objective is
$A+e^{-\lambda}[(y-1)A+\lambda(B-A)]$ with $B\ge A\ge0$.
Its derivative is
$e^{-\lambda}[(1-y)A+(1-\lambda)(B-A)]\ge0$ for $\lambda\le1$.
Increasing each coordinate below 1 therefore cannot decrease the value.
\end{proof}

\section{Small integers: retaining the probability-mass constraint}
For $1\le S\le n$, choose a category $h$ of maximal true probability,
before observing the sample. Its mean satisfies
\[
 \lambda_h=mp_h\ge m/S>\frac{70n}{SL}=:\Lambda.
\]
Hence
\begin{equation}\label{eq:anchor}
 \Prb(N_h<K)\le e^{-\Lambda}\sum_{\ell=0}^{K-1}\frac{\Lambda^\ell}{\ell!}
 \quad\hbox{for every positive integer }K.
\end{equation}
This step is essential: allowing arbitrary unconstrained means for every
category may give an upper bound above $1/3$ even when $S=1$.

Choose the exact envelope $q_j:=\max(0,\max_{\ell\ge3,\ \ell\ {
m odd}}u_{j,\ell})$.
Then
\[
 \sum_i(f_j(N_i)-1)
 \le q_j S+(u_{j,1}-q_j)F_1-(1+q_j)F_0.
\]
Consequently the scale's upper event implies
\begin{equation}\label{eq:F01}
 F_0<T_j(F_1),\qquad
 T_j(s)=\frac{(u_{j,1}-q_j)s+(q_j-\delta_j)S-\delta_jn}{1+q_j}.
\end{equation}
When $N_h\ge2$, $F_0,F_1$ come from the remaining $r=S-1$ categories.
Independence permits removing the anchor without changing their law.
Importantly, the $S$ in \eqref{eq:F01} remains the original support size.
For two rows, $n=16$ and $n=27$, use the stronger anchor event $N_h\ge4$.
In general, for an even $K\ge2$ put
\[
 q_{h,j}=\max(0,\max_{\ell\ge K+1,\ \ell\ {\rm odd}}u_{j,\ell}).
\]
On $N_h\ge K$ the anchor's correction is at most $q_{h,j}$. Replace
$T_j(s)$ in \eqref{eq:F01} by the sharper threshold
\[
 T_{j,K}(s)=T_j(s)-\frac{q_j-q_{h,j}}{1+q_j}.
\]
For $K=2$ this equals $T_j(s)$. This adjustment and the full Poisson
tail in \eqref{eq:anchor} are included in the certificates.

For a chosen collection of scales with $u_{j,1}\ge q_j$, define
$z_s=\lceil\max_j T_{j,K}(s)\rceil-1$.
This is the largest integer strictly below the threshold; certified
rational coefficient upper bounds may be used to enlarge it safely.
The sequence $z_s$ is nondecreasing. For any fixed rational $0<y<1$ put
\[
 \phi_y(s)=\begin{cases}0,&z_s<0,\\y^{-z_s},&z_s\ge0.\end{cases}
\]
Then
$1_{\{F_0\le z_{F_1}\}}\le y^{F_0}\phi_y(F_1)$.
The intersection of this union with $\{N_h\ge K\}$ is contained in the
remaining-coordinate event just bounded. Dropping the anchor condition
enlarges that event and does not condition the remaining variables.
Lemma~\ref{lem:homog} therefore bounds the intersection probability
by the supremum, over a single $\lambda\ge1$, of
\begin{equation}\label{eq:diagonal}
 \sum_{s=0}^r\binom rs
       (\lambda e^{-\lambda})^s
       [1-(1-y+\lambda)e^{-\lambda}]^{r-s}\phi_y(s).
\end{equation}
The remaining scales, if any, are bounded with \eqref{eq:chernoff} at
the specified $S$. Add these bounds and the anchor error
\eqref{eq:anchor}. The chosen split and rational $y,t,\eta$ may vary
with $n,S$: they are proof parameters, not changes to the estimator.
Empty scale collections and $r=0$ are handled directly.

The finite verifier includes every integer $1\le n\le73$ and
$1\le S\le n$, with complete scale coverage, exact strict-inequality
rounding, and monotonicity checks on $\phi$. It proves
\begin{equation}\label{eq:small}
 U_{1:73}<\SmallUpperDisplay.
\end{equation}
Here is the explicit enclosure of \eqref{eq:diagonal}. Write
$b(\lambda)=\lambda e^{-\lambda}$ and
$h(\lambda)=1-(1-y+\lambda)e^{-\lambda}$.
On $\lambda\ge1$, $b$ decreases and $h$ increases. On a closed interval
$[a,b_0]\subset[1,24]$, the positive sum is therefore at most
\[
 \sum_{s=0}^r\binom rs b(a)^s h(b_0)^{r-s}\phi_y(s).
\]
The verifier bisects until this outward interval bound is below the
specified rational target. For $\lambda\ge24$ use $h\le1$ and
$b(\lambda)\le24e^{-24}$ in the same positive sum.
The full certificate includes 2701 $(n,S)$ cases and covers every
$\lambda$ by these intervals and the explicit remaining half-line.
Its worst upper bound is below $0.330679396$, at $n=27$.
The bound for each fixed $S$ is assembled before taking the maximum;
unrelated maxima over $S$ are not multiplied or treated as independent.
Numerical optimization alone is never accepted as verification.

\section{Large range, coverage, and completion of the proof}
For $L\ge14$ the present estimator and the earlier $(70,14)$ estimator
are identical, including their real mean. The earlier concentration
proof, reproduced in the attached appendices, establishes joint failure
less than $0.0009$. It uses coefficient/range envelopes, Cantelli,
Hoeffding and endpoint Bernstein bounds, with 260 closed slabs covering
$[14,40]$ and a decreasing analytic envelope on $[40,\infty)$.
Thus its use here does not extrapolate a finite simulation to infinity.
For this proof the $C=70$ obligations have been freshly recomputed using
160-bit dyadic rational intervals: 15,885 checks, including equality with
the saved complete finite and half-line values. The final joint bound is
below $0.000879402$. The old counting branch is not used in this transfer.
In the combined PDF, Appendix A is the repository manuscript
\emph{A Single-Run Refinement of the Support-Testing Constants}
(\texttt{results\_parameter\_search/proof.pdf}), and Appendix B is
the earlier exact-certificate manuscript
(\texttt{artifacts/joint\_constants\_submission/submission\_proof.pdf}).
Their sources and arithmetic dependencies are included in the archive.
We use only their sufficiency results and supporting lemmas.

There is no gap between the new finite and continuous ranges because
$1+\log403<7<1+\log404$. The continuous interval includes its endpoints
and overlaps the large-range proof at 14. Thus the domain is exhausted:
\begin{center}
\begin{tabular}{lll}\toprule
Range & Upper deviation or joint bound & Coverage\\\midrule
$1\le n\le73$ & $U<\SmallUpperDisplay$ & all 2701 $(n,S)$ pairs\\
$n=74$ & $U<\BridgeUpperDisplay$ & all seven scales\\
$75\le n\le403$ & $U<0.291304$ & all 329 integers, 2570 scales\\
$7\le L\le14$ & $U\le10e^{-4}<0.183157$ & 700 closed slabs, 6597 branches\\
$L\ge14$ & joint $<0.0009$ & 260 closed slabs and an infinite tail\\\bottomrule
\end{tabular}
\end{center}
For $L<14$ and $S\le n$, $D\le S$ deterministically. Thus an upper
failure requires a polynomial scale event already bounded by $U$.
The lower interval failure is bounded by \eqref{eq:lower}, since
$\widetilde S\ge D$. Explicitly,
$0.330679396+1.657\cdot10^{-27}<0.331<1/3$; all other new upper
bounds are smaller. This rational margin is checked in the aggregate
certificate. On the null promise, the
same event also ensures the binary answer is correct.
If $S>n$, the clipped upper endpoint is automatic. On the alternative,
$D>n$ simultaneously gives the correct binary answer and interval;
its complement is bounded by the discovery lemma. In the gap only the
interval is required. These arguments also apply to countable support.
For $L\ge14$ use the transferred joint bound. This proves
Theorem~\ref{thm:main}.

\section{Verification, scope, and experimental interpretation}
The proof package separates exploratory optimizers from deterministic
certificate replay. Every accepted bound uses outward intervals and
global domain coverage. It includes independent coefficient and
polynomial-positivity implementations, missing-domain negative controls,
and hashes of the exact sources and certificate files. The single
entry point is \texttt{python -m proof\_l01.verify\_all} in the repository
environment. This is a computer-assisted mathematical proof, conditional
on the documented arithmetic software and analytic lemmas; it is not a
claim of proof-assistant formalization.

The theorem is about iid Poisson sampling from a fixed discrete
distribution. Existing real-data tests construct fixed empirical
distributions and resample from them, so they test precisely this
conditional model. They do not establish independent field observations,
known unseen ecological support, or a sample-efficiency advantage.
The observed 78 failures in 1.28 million primary paths, including a worst
cell rate of 1.1\%, are compatible with the theorem. The tested cells and
their confidence bounds remain empirical evidence separate from the
uniform proof. The cheaper certified count baseline and the lack of a
demonstrated practical sampling advantage remain unchanged conclusions.
The exact optimal real constant, attainment of that constant, and
necessity for every smaller constant are outside this sufficient-constant
result and are not needed for the validity of $C_0=70,L_0=1$.
\end{document}
